\section{问题一：预热平衡阶段}\label{sec:q1}

\subsection{求解}

问题一采用附录 2 的常物性，体积热容 $\rho c_p=\SI{2.132e6}{J/(m^3.K)}$，热扩散率 $\alpha=\SI{1.69e-7}{m^2/s}$，扩散系数只依赖含水率。此时式~\eqref{eq:govT} 是常系数线性方程，式~\eqref{eq:govC} 不含温度，温度与含水率互不影响。以 $N=1600$ 求解到 $\SI{1800}{s}$，表~\ref{tab:t1}、表~\ref{tab:t2} 给出指定时刻与位置的结果，完整结果见 result1.xlsx，径向分布见图~\ref{fig:q1}。

\begin{table}[H]\centering\small
\caption{30 分钟内药材的温度（\si{\celsius}）}\label{tab:t1}
\begin{tabular}{cccccc}
\toprule
 & \multicolumn{5}{c}{到药材中心的距离/cm} \\
\cmidrule(lr){2-6}
时间/s & 0 & 0.5 & 1 & 1.5 & 2 \\
\midrule
100 & 28.0001 & 28.0003 & 28.0040 & 28.0327 & 28.1801 \\
300 & 28.0408 & 28.0635 & 28.1514 & 28.3681 & 28.8487 \\
600 & 28.4533 & 28.5360 & 28.8039 & 29.3159 & 30.1652 \\
900 & 29.3243 & 29.4582 & 29.8755 & 30.6161 & 31.7304 \\
1200 & 30.5427 & 30.7098 & 31.2223 & 32.1126 & 33.4276 \\
1500 & 31.9957 & 32.1867 & 32.7660 & 33.7463 & 35.1205 \\
1800 & 33.5754 & 33.7720 & 34.3642 & 35.3621 & 36.7856 \\
\bottomrule
\end{tabular}

\end{table}

\begin{table}[H]\centering\small
\caption{30 分钟内药材的水分浓度（\si{kg/kg}）}\label{tab:t2}
\begin{tabular}{cccccc}
\toprule
 & \multicolumn{5}{c}{到药材中心的距离/cm} \\
\cmidrule(lr){2-6}
时间/s & 0 & 0.5 & 1 & 1.5 & 2 \\
\midrule
100 & 2.5500 & 2.5500 & 2.5500 & 2.5500 & 2.2470 \\
300 & 2.5500 & 2.5500 & 2.5500 & 2.5492 & 2.0508 \\
600 & 2.5500 & 2.5500 & 2.5500 & 2.5353 & 1.8770 \\
900 & 2.5500 & 2.5500 & 2.5497 & 2.5045 & 1.7547 \\
1200 & 2.5500 & 2.5500 & 2.5482 & 2.4646 & 1.6586 \\
1500 & 2.5500 & 2.5499 & 2.5445 & 2.4206 & 1.5787 \\
1800 & 2.5500 & 2.5497 & 2.5383 & 2.3755 & 1.5102 \\
\bottomrule
\end{tabular}

\end{table}

\begin{figure}[htbp]\centering
\includegraphics[width=\textwidth]{fig_q1.pdf}
\caption{问题一的温度与含水率径向分布（$100$--$\SI{1800}{s}$）。}
\label{fig:q1}
\end{figure}

\subsection{结果分析}

\parahead{温度由表及里升高，预热尚未完成} 烘房温度在 $\SI{1800}{s}$ 内由 $\SI{28}{\celsius}$ 升到 $\SI{41.5}{\celsius}$，药材由表面向内依次升温。$\SI{100}{s}$ 时只有表层开始升温（表面 $\SI{28.18}{\celsius}$，中心几乎不变）；$\SI{1800}{s}$ 时表面 $\SI{36.79}{\celsius}$、中心 $\SI{33.58}{\celsius}$，径向温差 $\SI{3.2}{K}$，中心仍比烘房低 $\SI{7.9}{K}$。此时 Fourier 数 $\alpha t/R_0^2=0.76$，热量刚传到中心，预热平衡远未完成。传热 Biot 数 $Bi_h=1.39$，表面换热与内部导热的阻力相当，所以既有表面与烘房的温差，也有明显的径向温差。

\parahead{失水集中在表层} 表面含水率在 $\SI{100}{s}$ 内就由 2.55 降到 2.2470，$\SI{1800}{s}$ 时降到 1.5102；而距中心 $\SI{1.5}{cm}$ 处降到 2.3755，$\SI{1}{cm}$ 处只降到 2.5383，$\SI{0.5}{cm}$ 以内基本不变。按含水率比初值低 0.01 计，失水层厚约 $\SI{1.0}{cm}$。水分扩散时间 $R_0^2/D\approx\SI{22.5}{h}$ 远大于 $\SI{1800}{s}$，且 $Bi_m=3.24$，表面失水快于内部补给，表层迅速变干，形成陡峭的含水率梯度。$\SI{1800}{s}$ 内平均含水率由 2.55 降到 2.294，干燥速率由 $\SI{0.729}{h^{-1}}$ 降到 $\SI{0.425}{h^{-1}}$。

预热平衡阶段的前 30 分钟里，热量已传到中心，失水却只发生在表层，传热与传质的快慢相差悬殊，这正是问题二中划分两个阶段的依据。
